\chapter{Methods}
\label{ch:methodsmain}

This chapter summarises what was designed, built and run. Full engineering
detail -- the master's electronics and firmware, the complete system
architecture, the perception and motion pipelines -- is carried in
\cref{ch:master,ch:method,ch:perception,ch:motion} so that it can be checked
and reproduced without lengthening the argument here.

\section{The master arm}

The master is a miniature mannequin reproducing the Gen3's joint sequence at
roughly a third of its scale (\cref{fig:master-cad}, \cref{ch:master}), moved
with the fingers of one or both hands. A Teensy 4.1 reads fourteen rotary
potentiometers, two inertial measurement units, two force-sensitive resistors
and two buttons, and emits one text frame at \SI{100}{\hertz}
(\cref{tab:frame}). Potentiometer counts pass through an integer exponential
moving average (\cref{eq:ema}) before conversion to degrees; the inertial
sensors are read as a single burst so that all six axes come from one
sampling instant.

Each channel carries a health verdict computed from its own recent history:
\emph{dead}, \emph{incoherent} (present but not tracing a trajectory),
\emph{intermittent}, or \emph{alive}. Validation is scoped to the channels
the active mode actually consumes, so a dead channel the mode does not read
cannot reject a frame, and when no well-conditioned channel remains the
system reports position \emph{unavailable} rather than extrapolating one.

\section{From master pose to arm command}

The commanded position is built in spherical form about the master's
shoulder rather than by forward kinematics from seven joint angles
(\cref{eq:elev,eq:azim,eq:reach}): elevation from the gravity vector,
azimuth from the shoulder roll, and radius from the forward-kinematic
magnitude of the bend joints. Each component is taken from whichever sensor
observes it most directly, so a dead roll channel can be passed through as
zero without corrupting the radius. The commanded end-effector position is
an anchor plus a scaled displacement from a reference latched at clutch
engagement (\cref{eq:anchor}), so releasing and re-engaging the clutch
indexes like a mouse lifted and set down again, and the reachable set is
bounded by the robot rather than by the master's own range.

\begin{table}[htbp]
  \centering
  \small
  \caption[Operating modes]{Operating modes, condensed from \cref{tab:modes}.
  Autonomy runs at a quarter of teleoperation's speed by design: nobody is
  holding a master, so the only bound on a wrong motion is how long it takes
  to happen.}
  \label{tab:modes-main}
  \begin{tabular}{@{}llll@{}}
    \toprule
    Mode & Input & Robot decides & $v_{\max}$ \\
    \midrule
    01/02 direct & mannequin or VR controllers & nothing & \SI{0.60}{\radian\per\second} \\
    03/04 shared autonomy & master or controllers, vision & wrist orientation & \SI{0.60}{\radian\per\second} \\
    06 full autonomy & typed or spoken goal & target, approach, grasp, place & \SI{0.15}{\radian\per\second} \\
    \bottomrule
  \end{tabular}
\end{table}

Five modes are delivered (\cref{tab:modes-main}); one safety stack serves all
of them, and dropping any shared mechanism is refused regardless of mode. The
desk arrangement the laboratory actually operates has the operator seated
across the room from the wearer, holding the VR controllers as a
six-degree-of-freedom input device and watching the physical robot directly
rather than wearing the headset (\cref{sec:desk}); the alignment between the
operator's frame and the robot's is a measured yaw rather than a mirror,
because facing someone and copying them is a reflection that would invert
every commanded orientation while positions looked correct.

\subsection{Safety architecture}

Collision avoidance is graduated in three stages (\cref{sec:clearance}):
collision-aware inverse kinematics, redundancy re-seeding on rejection, and a
geometric clearance floor of \SI{150}{\milli\metre} measured from the
transform tree between any moving arm link and any wearer primitive -- not
from the semantic collision description, which excludes the proximal pairs a
shoulder mount actually threatens. Every mechanism able to stop the arm
registers by name, must declare how it clears, and expires into an explicit
\emph{unknown} state (still treated as blocking) if the loop asserting it
stops running, because a blocker that silently keeps its last state is
indistinguishable from one that is working. Three independent triggers latch
the emergency stop until reset: an explicit call, both master buttons held
together, and a dead-man keyed on the \emph{source} timestamp of the master
data, not on message arrival -- the correction that makes it fire on the
failure mode that actually occurs, since a degraded master keeps publishing
stale values at full rate rather than falling silent.

\section{Perception, planning and motion}

A calibrate-then-plan stage (\cref{ch:perception}) replaced declared object
and surface coordinates with measured ones: a wrist-camera sweep that holds
one attitude per pass and translates, printer-bed-probe style, fits the
working surface and segments objects from the image rather than clustering a
point cloud, and hands the result to an RRT-Connect planner
\cite{kuffner2000rrtconnect} evaluated against the same moving-chain
clearance the safety layer enforces. Detection is open-vocabulary
(YOLO-World~\cite{cheng2024yoloworld} with FastSAM~\cite{zhao2023fastsam}),
falling back to a hue-and-depth filter, and grasp synthesis is the geometric
principal-axis method of \cref{eq:pca,eq:graspwidth} rather than a learned
one, chosen because it can \emph{name} the reason for a refusal and fits the
\SI{4}{\gibi\byte} of video memory available on this machine.

The per-joint step limiter that previously constituted the entire motion
generator for teleoperation left the commanded hand \SI{51.3}{\milli\metre}
off the path its own inverse-kinematics solution implied
(\cref{sec:clamp-maths}); it was replaced with Ruckig~\cite{berscheid2021ruckig},
a jerk-limited, phase-synchronised generator in which every joint is driven
from a single scalar phase (\cref{eq:phase-sync}), so the commanded
configuration is on the line between endpoints by construction. Replaying
recorded joint commands onto the physical arm shows every joint parks
\SI{0.305}{\degree} short of its target on the side it approached from
(\cref{sec:sim-to-real}); the two candidate models -- a terminal offset and a
proportional gain error -- are distinguishable only at short travel, and
propagating the error through the manipulator Jacobian selects the terminal
offset at $r=\num{0.923}$.

\section{Verification methodology}

Twenty-six times in this project a result that looked like a system failure
was a defect in the measuring tool (\cref{ch:instrument}). The practice that
follows from it is applied throughout: every analysis script carries a
known-answer test, synthetic data is used only where its ground truth is
\emph{constructed} rather than rendered, and a result contradicting an
earlier measurement is treated as an instrument check until the two are
reconciled. Because the inverse-kinematics solver restarts randomly, every
reachability claim in this report is ten repeats over a path densified to
\SI{20}{\milli\metre}, not single calls at the endpoints, and a pose that is
not solid on all ten repeats is not reported as reachable.

\section{Ethics and the pilot study}
\label{sec:ethics}

\textbf{This study was approved by Imperial College London's Science,
Engineering and Technology Research Ethics Committee, approval number
7111986.} All participants gave informed consent before taking part.
Participant data are anonymised throughout this report as P1--P4;
\repo{write\_manifest()} refuses to record a name, email address, date of
birth, address or telephone number, so anonymity is enforced in the code that
writes the session record rather than only in this text.

The full two-person protocol (\cref{app:study}) crosses direct teleoperation
against shared autonomy over five tasks with a separate wearer, and was
designed in full before this report but could not be run: it depends on
master channels that were incoherent (\cref{ch:results}), and on an approved
protocol that did not yet exist. Both blockers have since partly cleared.
Approval was obtained, and VR controllers -- already a delivered input mode
(\cref{tab:modes-main}) -- supply a six-degree-of-freedom command that does
not depend on the master's damaged channels, which made a first, smaller
pilot practical ahead of a full session with a separate wearer.

Four operators (P1--P4) each completed between two and four recorded
sessions, performing one of three manipulation-style tasks -- object
tracking, target reaching, and position matching -- once under direct VR
control and once with shared-autonomy assistance active, using the desk
arrangement of \cref{sec:desk}. Coverage across participant, task and
condition is uneven, which \cref{sec:pilot-results} reports rather than
hides. Every sample was logged automatically at \SI{20}{\hertz} -- controller
and end-effector pose, clutch engagement, and the VR link's own reported
spatial lag -- together with an event log of every clutch and emergency-stop
transition, giving completion time, hand-travel distance, re-grip count, the
active-control fraction of the session, and mean tracking lag per session
without any manual timing. This pilot measures \emph{operator} performance
only; it has no separate wearer and therefore does not address the
wearer-side hypotheses (H3, H4) of the full protocol, which remain for the
session \cref{app:study} specifies.
